% Part: sets-functions-relations
% Chapter: infinite
% Section: dedekinds-proof
%
\documentclass[../../../include/open-logic-section]{subfiles}

\begin{document}

\olfileid{sfr}{infinite}{dedekindsproof}

\olsection[د ډېډېکېنډ «ثبوت»]{د نامتناهي سټ د شتون په اړه د
ډېډېکېنډ «ثبوت»}

په دې څپرکي کښې مو د طبيعي عددونو سټي شننه د ډېډېکېنډ د الجبرونو
له مخې وړاندې کړه۔ په \olref[arith][ref]{sec} کښې مو (له نورو
خبرو سره) د تحليل د سټي جوړونې پر فلسفي اهميت غور وکړ۔ اوس بايد د
هغه څه اهميت وارزوو چې دلته مو ترلاسه کړل۔

د \olref[sfr][arith][]{chap} په ټوله موده کښې مو طبيعي عددونه
راکړل شوي وګڼل او په څرګند ډول مو ترې صحيح، ناطق او حقيقي عددونه
جوړ کړل۔ په دې څپرکي کښې مو د طبيعي عددونو کوم څرګند جوړښت نۀ دے
وړاندې کړے۔ يوازې مو وښودل چې \emph{هر د ډېډېکېنډ له مخې نامتناهي
سټ راکړل شي}، داسې سټ تعريفولے شو چې کټ مټ لکه څنګه چې غواړو
$\Nat$ وچلېږي، هماغسې به وچلېږي۔

نو څرګنده ده چې دا ادعا نشو کولے چې کومې ماوراءالطبيعتي پوښتنې ته
مو ځواب ورکړے، لکه دا چې \emph{طبيعي عددونه کوم څيزونه دي}۔ خو دا
ښۀ کار دے۔ په \olref[sfr][arith][ref]{sec} کښې مو ټينګار کړے و چې
د ډېډېکېنډ د پرېکونونو د سټ په توګه د $\Real$ تعريف بايد يوه اسانه
ټاکنه وګڼو، نۀ يوه \emph{موندنه}۔ مهمه مشاهده دا ده چې د
ډېډېکېنډ پرېکونونه د حقيقي عددونو بنسټيز رياضي خاصيتونه لري۔ دلته
هم همدا خبره ده: مهمه مشاهده دا ده چې \emph{هر} ډېډېکېنډ الجبر د
طبيعي عددونو بنسټيز رياضي خاصيتونه لري۔ (په رښتيا، ډېډېکېنډ دا
ټکے په دې ثبوت لا پياوړے کړ چې ټول ډېډېکېنډ الجبرونه
\emph{آيزومورف} دي (\citeyear[Theorems
132--3]{Dedekind1888})۔ نو
دا هېڅ د حيرانۍ خبره نۀ ده چې اوسني ډېر «جوړښت‌پال» ډېډېکېنډ خپل
مخکښ بولي۔)

برسېره پر دې، موږ وښودل چې د طبيعي عددونو نظريه څنګه په ساده سټ
نظريه کښې ورګډېږي؛ دا ساده نظريه خپله لا هم يو څه غېر رسمي ده، خو
(لکه چې ښکاري) طبيعي عددونه له مخکښې راکړل شوي نۀ ګڼي۔ نو کېدے شي
د ډېډېکېنډ د خپلې سترې پروژې د پوره کولو پر لار يو، چې هغه داسې
تشريح کړې وه:
\begin{quote}
	په ساينس کښې هېڅ د ثبوت وړ خبره بايد بې له ثبوته ونۀ منل شي۔ که
	څه هم دا غوښتنه معقوله ښکاري، زۀ دا نشم منلے چې آن د تر ټولو ساده
	ساينس د بنسټ ايښودلو په وروستيو طريقو کښې پوره شوې ده؛ يعنې د منطق
	په هغې برخې کښې چې د عددونو نظريه څېړي۔ کله چې حساب (الجبر،
	تحليل) يوازې د منطق يوه برخه بولم، مقصد مې دا دے چې د عدد مفهوم
	د مکان او زمان له تصورونو يا شهود څخه بشپړ خپلواک ګڼم---بلکې دا
	د فکر د خالصو قوانينو نېغه پايله ګڼم۔
	\citep[preface]{Dedekind1888}
\end{quote}
د ډېډېکېنډ زړه‌وره مفکوره دا ده۔ همدا اوس مو وښودل چې يوازې د
(ساده) سټ نظريې په کارولو طبيعي عددونه څنګه جوړېږي۔ په
\olref[sfr][arith][]{chap} کښې مو وليدل چې د طبيعي عددونو او يوې
اندازې سټ نظريې په راکړه کېدو سره حقيقي عددونه څنګه جوړېږي۔ نو
ښايي «حساب (الجبر، تحليل)» «يوازې د منطق يوه برخه» وګرځي (د
ډېډېکېنډ د «منطق» کلمې په پراخه معنا)۔

نظر همدا دے۔ خو يوه شېبه تم شئ۔ زموږ د ډېډېکېنډ الجبر جوړول (چې
د طبيعي عددونو بدل دے) د يوه د ډېډېکېنډ له مخې نامتناهي سټ د شتون
پر شرط ولاړ دي۔ (يوازې \olref[sfr][infinite][dedekind]{thm:DedekindInfiniteAlgebra}
ته بېرته وګورئ۔) تر څو چې د ډېډېکېنډ له مخې د نامتناهي سټ شتون د
«منطق» يا «د فکر د خالصو قوانينو» له لارې ثابت نشي، پروژه درېږي۔

نو، \emph{ايا} د ډېډېکېنډ له مخې د نامتناهي سټ شتون «د فکر د
خالصو قوانينو» له لارې ثابتېدے شي؟ د ډېډېکېنډ هڅه دا وه:
\begin{quote}
  زما د فکرونو خپله نړۍ، يعنې د ټولو هغو څيزونو بشپړتيا $S$ چې زما
  د فکر موضوع کېدے شي، نامتناهي ده۔ ځکه که $s$ د~$S$ يو غړے وي،
  نو دا فکر $s'$ چې~$s$ زما د فکر موضوع کېدے شي، خپله د~$S$ يو
  غړے دے۔ که دا د غړي~$s$ يو تصوير $\phi(s)$ وګڼو، نو
  \dots~$S$ [د ډېډېکېنډ له مخې] نامتناهي دے، او همدا بايد ثابت
  شوي وے۔
	\citep[\S66]{Dedekind1888}
\end{quote}
داسې خبره د يوه داسې کتاب په منځ کښې ليدل رښتيا هم ډېره حيرانوونکې
ده چې زياته برخه ئې له بشپړو دقيقو رياضي ثبوتونو جوړه ده۔ دوه
تبصرې د پام وړ دي۔

لومړۍ: دا «ثبوت» هغه بڼه نږدې هېڅ نۀ لري چې اوس به ئې «رياضي»
وبولو۔ دا د نفسياتي څيزونو (فکرونو) خبره کوي، او هغه هم يوازې
\emph{ممکن} څيزونه۔

دويمه: لږ تر لږه لکه څنګه چې موږ د ډېډېکېنډ الجبرونه وړاندې کړي،
دا «ثبوت» يوه څرګنده فني نيمګړتيا لري۔ که د ډېډېکېنډ استدلال بريالے
وي، يوازې دا ثابتوي چې نامتناهي ډېر څيزونه (په ځانګړې توګه
نامتناهي ډېر فکرونه) شته۔ خو ډېډېکېنډ بايد دا دليل هم راکړي چې~$S$
ولې يو يوازينے \emph{سټ} وګڼو چې نامتناهي ډېر !!{element}s لري،
نۀ دا چې~$S$ \emph{څه څيزونه} (په جمع بڼه) وګڼو۔

دا حقيقت چې ډېډېکېنډ دلته تشه نۀ ليدله، ښايي وښيي چې د هغه د
«بشپړتيا» کلمه کټ مټ زموږ د «سټ» کلمې له کارونې سره برابره نۀ
ده۔\footnote{په رښتيا، د دې د باور دپاره نور دليلونه هم لرو؛
\citet[p.~23]{Potter2004} وګورئ۔} خو دا به ډېره د حيرانۍ خبره نۀ
وي۔ هغه پروژه چې په تېرو دوو څپرکيو کښې مو مخ ته يوړه---د طبيعي
عددونو «جوړول»، او له هغو څخه د صحيح، حقيقي او ناطقو عددونو
«جوړول»---ټوله په ساده ډول شوې ده۔ کله مو چې کوم سټ ته اړتيا لرله،
هماغه سټ مو منلے؛ خو د دې په اړه مو ډېر عمومي اصلونه نۀ دي ټاکلي
چې کټ مټ کوم سټونه کله شتون لري۔ خو پوهېږو چې \emph{ځينې} عمومي
اصلونه پکار دي، ځکه له هغو پرته به د رسل په پاراډوکس کښې ولوېږو۔

اوس د دې وخت دے چې له خپلې ساده‌ګۍ څخه ووځو۔

\end{document}
